\section{Discussion}
\subsection{Interpretation and Answers to the Research Questions}
For RQ1, prediction-triggered reservation improved restricted mission time relative to the reactive short queue, but the contrast is dominated by an identified failure mode: boats can become idle beyond radio range. It does not isolate forecast skill. For RQ2, the always-long queue performed at least as well in mean restricted time in every condition, and fixed partitioning completed reliably with less coordination cost. Thus these experiments do not establish that the predictive policy is preferable. For RQ3, held-out replay showed similar, low event-level F1 for trend and persistence. Abrupt Markov fade onsets are not necessarily predictable from received history.

The revised plan anticipated both null effects and stalled reconnection. Retaining those outcomes exposes a limitation of the chosen coordination model. Increasing queue size can prevent some queue exhaustion even without a useful forecast. The present results support that mechanism more directly than a prediction benefit. Redundant observations arise from transit and do not imply inconsistent ownership.

\subsection{Validity, Limitations, and Generalizability}
The evidence concerns one grid, two start locations, one common greedy allocator, and synthetic numerical controls. S2's complete censoring makes its restricted mean well-defined but leaves its uncensored completion-time distribution unidentified. Percentile bootstrap intervals with 30 seeds can be discrete and sensitive to a few stalled missions; the six conditions should not be pooled into a universal effect. Reusing the same fade trace across policies controls exogenous randomness but cannot equalize route-dependent delivery.

Atomic batch commitment simplifies agreement and omits partial delivery, retransmission timing, contention, and real distributed consensus. Ideal cell-centre sensing omits sonar, localization error, and bathymetric uncertainty. The simulator has no rendezvous or collision avoidance, so it is unsuitable as an operational controller. All results require independent author reproduction and review. Software verification does not calibrate a marine environment.

\subsection{Ethics, Societal Aspects, and Sustainability}
Synthetic data avoid human-participant privacy issues. The main responsibility is to prevent a small verified model from being presented as evidence of safe autonomous vessel deployment, accurate seabed maps, or measured energy savings. A stalled policy can appear cheap in movement or communication because it accomplishes less work; coverage and censoring must accompany resource indicators. Bytes, operation counts, laptop CPU time, and process RAM do not determine radio or propulsion energy, battery endurance, or device lifetime.

\subsection{Implications and Future Work}
A separate, explicitly revised study could add one common recovery/rendezvous rule and repeat all controls, or test a channel with independently justified observable precursors. These changes should be specified before a new evaluation rather than tuned to reverse the current ranking. Measured radio traces and vessel/sensor models would then be needed to assess external validity. The existing long-queue ablation should remain, because it reveals when reservation capacity accounts for the apparent benefit of prediction.
